\documentclass[11pt,a4paper]{article}

\usepackage{fontspec}
\usepackage{polyglossia}
\setmainlanguage{russian}
\setotherlanguage{english}

\IfFontExistsTF{Times New Roman}
{\setmainfont{Times New Roman}}
{\setmainfont{Libertinus Serif}}

\usepackage[a4paper,top=1.8cm,bottom=1.8cm,left=1.9cm,right=1.9cm]{geometry}
\usepackage{xcolor}
\usepackage{titlesec}
\usepackage{enumitem}
\usepackage{amsmath,amssymb}
\usepackage{booktabs}
\usepackage{longtable}
\usepackage{hyperref}

\definecolor{hseblue}{HTML}{144F7B}
\definecolor{softblue}{HTML}{EAF4FA}
\definecolor{darkred}{HTML}{8B1E3F}

\hypersetup{colorlinks=true, linkcolor=black, urlcolor=hseblue}

\setlength{\parindent}{0pt}
\setlength{\parskip}{0.5em}
\linespread{1.05}\selectfont
\setlist[itemize]{leftmargin=1.5em, topsep=2pt, itemsep=2pt}
\setlist[enumerate]{leftmargin=1.6em, topsep=2pt, itemsep=2pt}

\titleformat{\section}
  {\Large\bfseries\color{hseblue}}
  {\thesection.}
  {0.5em}
  {}

\titleformat{\subsection}
  {\large\bfseries\color{hseblue}}
  {\thesubsection.}
  {0.5em}
  {}

\titleformat{\subsubsection}
  {\bfseries\color{darkred}}
  {}
  {0pt}
  {}

\newcommand{\answerbox}[1]{%
  \noindent\colorbox{softblue}{%
    \parbox{\dimexpr\linewidth-2\fboxsep\relax}{\textbf{Короткий ответ:} #1}%
  }\par\vspace{0.2em}
}

\begin{document}

\begin{center}
{\Large\bfseries Подробные ответы к выступлению\par}
\vspace{0.35em}
{\large Математика, смысл метрик, тесты и возможные вопросы комиссии\par}
\vspace{0.45em}
Ковалёва Мария Сергеевна \\
НИУ ВШЭ, Факультет экономических наук, Москва, 2026
\end{center}

\vspace{0.5em}
Этот файл дополняет краткий текст выступления. Его задача --- быстро закрыть вопросы по методологии, формулам, интерпретации метрик, статистическим тестам и ограничениям работы.

\tableofcontents
\newpage

\section{Карта защиты по слайдам}

\begin{longtable}{p{0.16\linewidth}p{0.39\linewidth}p{0.36\linewidth}}
\toprule
\textbf{Слайд} & \textbf{Что доказывает} & \textbf{Главный ответ на вопрос} \\
\midrule
1. Тема & Работа эмпирическая: сравниваются модельные и рыночные цены. & Исследуется не «истинная цена», а устойчивость ошибки модели на данных MOEX. \\
2. Проблема & Модели используют мало параметров, рынок агрегирует больше информации. & Ошибка ожидаема; важен её масштаб и структура. \\
3. Гипотеза & Ошибка зависит от режима рынка и может расти в стрессовых состояниях. & Гипотеза подтверждается в уточнённой форме: режим важен, но DTE и денежность объясняют большую часть эффекта. \\
4. Методология & Данные $\rightarrow$ волатильность $\rightarrow$ модель $\rightarrow$ ошибка $\rightarrow$ тесты. & Все модели сравниваются на общей выборке, чтобы различие не было следствием разных наблюдений. \\
5. Модели & CRR учитывает американское исполнение, Black-76 нет. & Для американских опционов на фьючерс CRR методологически корректнее. \\
6. Общее сравнение & GARCH и HV\_63d почти равны по MAE; все модели недооценивают рынок. & Простая модель с одной $\sigma$ не устраняет систематическое смещение. \\
7. Режимы & GARCH лучше в low/mid vol, хуже в high vol. & В стрессе GARCH может быть инерционным, а состав выборки меняется. \\
8. TVNAE & GARCH лучше описывает временную стоимость. & TVNAE выделяет именно ту часть цены, которую должна объяснять модель. \\
9. CRR vs Black-76 & Black-76 сильно хуже, есть премия раннего исполнения. & Право досрочного исполнения имеет экономическую стоимость. \\
10. Тесты & Смещение значимо, GARCH и HV\_63d статистически близки, HV\_21d хуже. & Выводы не только описательные, они проверены формально. \\
11. Выводы & Сводятся пять главных результатов. & На коротких горизонтах полезен GARCH, на длинных устойчив HV\_63d, всем моделям не хватает поверхности волатильности. \\
12. Ограничения & Период, рынок, SETTLEPRICE, одна $\sigma$, постоянная ставка. & Ограничения не ломают вывод, но задают направления развития. \\
13. Финал & Готовность отвечать. & Самый короткий итог: американская постановка важна, GARCH полезен для временной стоимости, простые модели занижают рынок. \\
\bottomrule
\end{longtable}

\section{Данные и переменные}

\subsection{Что именно является рыночной ценой}

\answerbox{Рыночная цена в работе --- расчётная цена Московской биржи, SETTLEPRICE. Это воспроизводимый биржевой бенчмарк, а не утверждение об абсолютно истинной фундаментальной стоимости.}

В работе используются дневные наблюдения по опционам серии MM на фьючерс MXI за период с 3 января 2024 года по 19 мая 2026 года. MXI --- фьючерс на Индекс МосБиржи. Это важно: в моделях для фьючерсных опционов базовым активом является не сам индекс, а цена фьючерса.

Исходная база содержит 19 574 строки, 92 уникальных опционных инструмента и 5 базовых фьючерсных контрактов. После фильтрации некорректных наблюдений остаётся 19 565 строк. Для основного сравнения трёх американских моделей используется общая выборка из 18 058 наблюдений, где одновременно доступны все три оценки волатильности.

\subsection{Основные обозначения}

\begin{itemize}
    \item $F$ --- расчётная цена базового фьючерса.
    \item $K$ --- страйк опциона.
    \item $T$ --- время до экспирации в годах.
    \item DTE --- срок до экспирации в календарных днях.
    \item $r$ --- безрисковая ставка, в базовой спецификации 16\% годовых.
    \item $\sigma$ --- годовая волатильность, заданная одним из трёх способов.
    \item $\widehat{P}$ --- модельная цена опциона.
    \item $P$ --- рыночная расчётная цена опциона.
\end{itemize}

\subsection{Денежность}

Денежность показывает, где находится цена базового фьючерса относительно страйка:

\[
M=\frac{F}{K}, \qquad \ln M=\ln(F/K).
\]

Для колла $F>K$ означает «в деньгах», для пута наоборот. В регрессии используется абсолютная логарифмическая денежность $|\ln(F/K)|$, потому что она показывает удалённость опциона от состояния около денег независимо от направления.

\subsection{Внутренняя и временная стоимость}

Внутренняя стоимость --- это то, что можно получить при немедленном исполнении:

\[
H_{\text{call}}(F,K)=\max(F-K,0), \qquad
H_{\text{put}}(F,K)=\max(K-F,0).
\]

Временная стоимость --- остаток рыночной цены сверх внутренней:

\[
TV=P-H(F,K).
\]

Экономический смысл: внутренняя стоимость почти механическая, а временная стоимость отражает неопределённость будущего движения, срок до экспирации, ставку и волатильность. Поэтому для оценки качества волатильностной модели особенно важно смотреть именно на временную стоимость.

\section{Оценка волатильности}

\subsection{Почему волатильность центральна}

\answerbox{Волатильность --- главный параметр неопределённости. Чем выше ожидаемая волатильность, тем дороже опцион, потому что у держателя больше шанс получить выгодный payoff.}

Опцион имеет асимметричный payoff: потери держателя ограничены премией, а потенциальный выигрыш зависит от движения базового актива. Поэтому рост неопределённости обычно увеличивает цену как колла, так и пута.

\subsection{Дневные лог-доходности}

Волатильность оценивается по дневным логарифмическим доходностям фьючерса:

\[
r_t=\ln\left(\frac{F_t}{F_{t-1}}\right).
\]

Лог-доходности удобны тем, что лучше агрегируются во времени и часто используются в моделях финансовых рядов.

\subsection{Историческая волатильность HV\_21d и HV\_63d}

Историческая волатильность --- стандартное отклонение дневных доходностей на окне, приведённое к годовому масштабу:

\[
\sigma_{21,t}=\sqrt{252}\,\operatorname{sd}(r_{t-20},\ldots,r_t),
\]

\[
\sigma_{63,t}=\sqrt{252}\,\operatorname{sd}(r_{t-62},\ldots,r_t).
\]

Множитель $\sqrt{252}$ используется, потому что в году примерно 252 торговых дня. 21 день --- около одного торгового месяца, 63 дня --- около квартала.

HV\_21d быстрее реагирует на свежие движения, но более шумная. HV\_63d устойчивее и лучше соответствует выборке, где преобладают длинные опционы.

\subsection{GARCH(1,1)}

GARCH моделирует условную дисперсию как зависящую от прошлых шоков и прошлой дисперсии:

\[
r_t=\varepsilon_t,\qquad \varepsilon_t=\sigma_t z_t,
\]

\[
\sigma_t^2=\omega+\alpha\varepsilon_{t-1}^2+\beta\sigma_{t-1}^2.
\]

Здесь $\omega$ --- базовый уровень дисперсии, $\alpha$ показывает реакцию на новый шок, $\beta$ показывает персистентность волатильности. Если $\alpha+\beta$ близко к 1, шоки затухают медленно.

В работе GARCH оценивается на расширяющемся окне: модель использует только информацию, доступную на момент прогноза. Это важно для защиты от look-ahead bias, то есть от использования будущих данных.

\subsection{Почему GARCH выигрывает на коротких DTE}

GARCH даёт условный прогноз текущей волатильности. Это особенно полезно для коротких опционов, цена которых сильно зависит от ближайшей турбулентности. Для длинных опционов важна ожидаемая волатильность на месяцы вперёд, и однодневный прогноз GARCH уже не отражает весь горизонт.

\subsection{Почему GARCH хуже в high vol}

В стрессовом режиме GARCH может переоценивать устойчивость текущего шока или, наоборот, запаздывать при резкой смене режима. Кроме того, в high-vol подвыборке меняется состав контрактов: больше длинных и специфических по денежности наблюдений. Поэтому сам режим не является единственной причиной ошибки.

\section{Модели ценообразования}

\subsection{CRR American}

\answerbox{CRR --- биномиальная модель, которая строит дерево возможных цен и на каждом узле выбирает максимум между немедленным исполнением и продолжением. Поэтому она подходит для американских опционов.}

В модели CRR на каждом шаге цена фьючерса может пойти вверх или вниз:

\[
u=e^{\sigma\sqrt{\Delta t}}, \qquad d=e^{-\sigma\sqrt{\Delta t}}=\frac{1}{u}.
\]

Для опциона на фьючерс риск-нейтральная вероятность в дереве задаётся так:

\[
q=\frac{1-d}{u-d}.
\]

Интуиция: в риск-нейтральном мире фьючерсная цена является мартингалом, поэтому ожидаемая будущая фьючерсная цена равна текущей. Для акции в классической CRR-формуле появляется $e^{r\Delta t}$, а для фьючерса центрирование другое.

На последнем шаге считается payoff:

\[
C_T=\max(F_T-K,0), \qquad P_T=\max(K-F_T,0).
\]

Затем стоимость возвращается назад по дереву. Для американского опциона на каждом узле:

\[
V=\max\left(\text{intrinsic value},\ e^{-r\Delta t}\left[qV_u+(1-q)V_d\right]\right).
\]

Первый аргумент --- немедленное исполнение, второй --- стоимость продолжения. Именно это делает модель американской.

\subsection{Black-76}

\answerbox{Black-76 --- европейская аналитическая формула для опционов на фьючерсы. Она удобна как benchmark, но не учитывает досрочное исполнение.}

Формулы:

\[
C=e^{-rT}\left(FN(d_1)-KN(d_2)\right),
\]

\[
P=e^{-rT}\left(KN(-d_2)-FN(-d_1)\right),
\]

\[
d_1=\frac{\ln(F/K)+0.5\sigma^2T}{\sigma\sqrt{T}}, \qquad
d_2=d_1-\sigma\sqrt{T}.
\]

Здесь $N(\cdot)$ --- функция распределения стандартной нормальной величины. Важное отличие от Black--Scholes: в Black-76 базовым активом является фьючерсная цена $F$, поэтому формула естественно подходит для фьючерсных опционов.

\subsection{Почему Black-76 хуже для этой выборки}

Исследуемые опционы допускают американское исполнение. Black-76 предполагает исполнение только на экспирации. Поэтому она не может учесть стоимость права исполнить опцион раньше. В работе MAE у Black-76 + HV\_63d равен 156,2 пункта против 107,1 у CRR + HV\_63d. Разница примерно 46\%.

Особенно сильный аргумент: в 24,5\% наблюдений цена Black-76 оказалась ниже внутренней стоимости. Для американского опциона это экономически некорректно, потому что держатель мог бы исполнить опцион немедленно.

\subsection{Премия раннего исполнения}

Премия раннего исполнения считается как:

\[
\Pi^R=V^{CRR}_{63}-V^{Black}_{63}.
\]

Средняя премия в работе --- 59,28 пункта, медиана --- 9,87. На коротких сроках она почти нулевая, а на длинных контрактах резко растёт. Это логично: чем больше времени до экспирации и чем выше ставка, тем ценнее возможность управлять моментом исполнения.

\section{Метрики ошибок}

\subsection{Знаковая ошибка}

\[
e_i=\widehat{P}_i-P_i.
\]

Если $e_i>0$, модель переоценивает опцион. Если $e_i<0$, модель недооценивает опцион. В работе средняя ошибка отрицательна у всех моделей, то есть все они систематически занижают биржевые расчётные цены.

\subsection{MAE}

\[
MAE=\frac{1}{n}\sum_{i=1}^{n}|e_i|.
\]

MAE показывает средний размер ошибки в пунктах цены. Это самая понятная метрика: «на сколько пунктов в среднем модель промахивается».

\subsection{RMSE}

\[
RMSE=\sqrt{\frac{1}{n}\sum_{i=1}^{n}e_i^2}.
\]

RMSE сильнее штрафует крупные ошибки, потому что ошибка сначала возводится в квадрат. Если у модели MAE близкий, но RMSE выше, значит у неё больше крупных промахов.

\subsection{Mean error}

\[
\overline{e}=\frac{1}{n}\sum_{i=1}^{n}e_i.
\]

Mean error показывает направление смещения. В работе:

\begin{itemize}
    \item CRR + GARCH: $-78{,}1$ пункта;
    \item CRR + HV\_63d: $-89{,}7$ пункта;
    \item CRR + HV\_21d: $-95{,}3$ пункта.
\end{itemize}

Отрицательный знак означает, что модельные цены ниже рыночных.

\subsection{TVNAE}

\answerbox{TVNAE --- это абсолютная ошибка, делённая на временную стоимость опциона. Метрика показывает, насколько модель ошибается относительно той части цены, которую действительно должна объяснять.}

\[
TVNAE_i=\frac{|\widehat{P}_i-P_i|}{TV_i}, \qquad TV_i=P_i-H(F_i,K_i).
\]

Если TVNAE = 0,20, ошибка равна 20\% временной стоимости. Если TVNAE = 1, модель ошиблась на величину всей временной стоимости.

TVNAE особенно нужна, потому что выборка смещена в сторону длинных опционов в деньгах. У таких контрактов большая часть цены --- внутренняя стоимость. MAE может выглядеть большой просто из-за масштаба цены, а TVNAE показывает качество именно опциональной компоненты.

\subsection{Почему не хватило MAPE}

MAPE делит ошибку на полную цену опциона:

\[
MAPE=\frac{1}{n}\sum_{i=1}^{n}\left|\frac{e_i}{P_i}\right|.
\]

Но полная цена включает внутреннюю стоимость. Для deep ITM опциона модель может хорошо воспроизвести внутреннюю стоимость и плохо оценить временную, а MAPE это сгладит. TVNAE решает именно эту проблему.

\subsection{Почему TVNAE считается только при $TV\geq10$}

Если временная стоимость близка к нулю, знаменатель очень мал, и метрика становится нестабильной: даже крошечная абсолютная ошибка даёт огромное значение. Порог 10 пунктов убирает такие технически нестабильные случаи. После порога остаётся 16 961 наблюдение, то есть 93,9\% общей выборки.

\section{Статистические тесты}

\subsection{Тест средней ошибки}

\answerbox{Тест проверяет, отличается ли средняя ошибка модели от нуля. В работе она значимо отрицательная для всех моделей.}

Регрессионная форма:

\[
e_{i,t}=\mu+\varepsilon_{i,t}.
\]

Проверяется гипотеза:

\[
H_0:\mu=0.
\]

Если $p<0{,}05$, нулевая гипотеза отвергается. В работе $p<0{,}001$ для всех трёх моделей. Значит, систематическая недооценка статистически значима.

Стандартные ошибки кластеризуются по торговой дате, потому что ошибки разных контрактов в один день не независимы: на них действует общий рыночный фон.

\subsection{Diebold--Mariano}

\answerbox{DM-test сравнивает прогнозное качество двух моделей по функции потерь. Он отвечает не «какая MAE меньше в таблице», а статистически значима ли разница качества.}

Для пары моделей $A$ и $B$ строится дневная разность потерь:

\[
d_t=\overline{L}^A_t-\overline{L}^B_t.
\]

Если используется абсолютная потеря, $L=|e|$. Если квадратичная, $L=e^2$. Нулевая гипотеза:

\[
H_0:\mathbb{E}[d_t]=0.
\]

В работе разность проверяется регрессией:

\[
d_t=\delta+u_t.
\]

Используются HAC-ошибки Newey--West, потому что дневные разности могут быть автокоррелированы: качество модели сегодня связано с качеством вчера.

Главные результаты:

\begin{itemize}
    \item GARCH и HV\_63d статистически неразличимы на полной выборке: $p=0{,}770$ по MAE и $p=0{,}370$ по MSE.
    \item HV\_21d статистически значимо хуже HV\_63d: $p=0{,}0003$.
\end{itemize}

\subsection{OLS-регрессия абсолютной ошибки}

\answerbox{OLS здесь не модель цены, а описательный инструмент: он показывает, какие характеристики связаны с масштабом ошибки.}

Финальная спецификация:

\[
|e_i|=\beta_0+\beta_1\mathbf{1}[\text{high\_vol}]_i+\beta_2DTE_i+\beta_3|\ln(F_i/K_i)|+\beta_4\mathbf{1}[\text{put}]_i+\varepsilon_i.
\]

Смысл переменных:

\begin{itemize}
    \item $\mathbf{1}[\text{high\_vol}]$ --- индикатор режима высокой волатильности.
    \item $DTE$ --- срок до экспирации в днях.
    \item $|\ln(F/K)|$ --- удалённость от состояния около денег.
    \item $\mathbf{1}[\text{put}]$ --- индикатор пута.
\end{itemize}

В финальной спецификации $R^2\approx0{,}27$. Это означает, что выбранные характеристики объясняют около 27\% вариации абсолютной ошибки. Для рыночных данных это содержательный результат: большая часть ошибки всё равно остаётся связанной с факторами вне простой модели, например ликвидностью и подразумеваемой волатильностью.

Ключевая интерпретация: при добавлении DTE и денежности индикатор high vol меняет знак. Значит, простой рост ошибки в высоковолатильном режиме во многом связан с составом наблюдений, а не с режимом как самостоятельной причиной.

\section{Как объяснять основные результаты}

\subsection{Почему все модели недооценивают рынок}

Главные причины:

\begin{itemize}
    \item используется одна плоская оценка волатильности, а реальный рынок имеет улыбку и срочную структуру implied volatility;
    \item расчётная цена биржи может включать ликвидностную премию и особенности биржевой методологии;
    \item ставка задана константой, хотя реальная срочная структура ставок неоднородна;
    \item модели не учитывают jumps, стохастическую волатильность, спрос и предложение, bid-ask spread и открытый интерес.
\end{itemize}

Главная формулировка для защиты: «Систематическая недооценка показывает не ошибку вычисления, а ограничение простой постановки с одной волатильностью на контракт».

\subsection{Почему HV\_21d хуже}

21 торговый день --- короткое окно. Оно быстро реагирует на свежие шоки, но даёт шумную оценку. Для выборки, где много длинных опционов, такая краткосрочная волатильность плохо соответствует горизонту, на котором формируется цена. Поэтому HV\_21d хуже по MAE и статистически уступает HV\_63d.

\subsection{Почему HV\_63d хорошая baseline-модель}

63 дня --- более устойчивое квартальное окно. Оно сглаживает случайные всплески и лучше подходит для длинных контрактов, доминирующих в выборке. На полной выборке HV\_63d почти равна GARCH по MAE и лучше по RMSE, то есть лучше контролирует крупные промахи.

\subsection{Почему GARCH полезен}

GARCH адаптируется к текущей условной волатильности. Поэтому он особенно полезен для коротких и среднесрочных опционов, а также по TVNAE, где оценивается именно временная стоимость. В работе GARCH имеет TVNAE 0,66 против 0,70 у HV\_63d, а на DTE 0--7 дней преимущество особенно сильное.

\subsection{Почему класс модели важнее оценки волатильности}

Разница между CRR и Black-76 больше, чем различия между HV\_21d, HV\_63d и GARCH внутри CRR. Это означает, что сначала нужно правильно выбрать тип модели под контракт. Для американского опциона модель должна учитывать право досрочного исполнения.

\section{Готовые ответы на вероятные вопросы}

\subsubsection{Почему не использовалась подразумеваемая волатильность?}

Подразумеваемая волатильность фактически извлекается из рыночной цены самого опциона. Если использовать её для проверки способности модели воспроизвести рыночную цену, возникнет элемент круговой логики. В этой работе задача была сравнить простые воспроизводимые оценки волатильности: исторические окна и GARCH. При этом IV surface указана как главное направление развития.

\subsubsection{Почему рыночная цена --- SETTLEPRICE, а не сделки?}

SETTLEPRICE есть почти для всех наблюдений и даёт сопоставимый дневной бенчмарк. Сделочные цены и bid-ask лучше отражали бы ликвидность, но по многим опционам могут быть редкими или шумными. Поэтому SETTLEPRICE выбран для воспроизводимости, а ликвидность вынесена в ограничения.

\subsubsection{Почему ставка фиксирована на 16\%?}

16\% используется как базовая постановка, соответствующая выбранному уровню ключевой ставки в расчётах. Чтобы проверить, не держится ли результат только на этой ставке, сделана чувствительность на 12\%, 16\% и 20\%. MAE меняется примерно на несколько пунктов, но ранжирование моделей качественно сохраняется.

\subsubsection{Почему используется 100 шагов CRR?}

100 шагов дают баланс между точностью и вычислительной устойчивостью на большой выборке. Биномиальная модель сходится к непрерывному пределу при росте числа шагов, но увеличение шагов резко повышает вычислительную нагрузку. Для целей сравнительного анализа 100 шагов достаточно как стабильная практическая спецификация.

\subsubsection{Что такое риск-нейтральная вероятность простыми словами?}

Это техническая вероятность, при которой дисконтированная ожидаемая цена соответствует текущей цене. Она не обязана совпадать с реальной вероятностью роста или падения. Её смысл --- корректно оценить цену производного инструмента без арбитража.

\subsubsection{Почему в Black-76 нет ожидаемой доходности актива?}

Потому что в риск-нейтральном ценообразовании ожидаемая доходность заменяется безрисковой логикой дисконтирования. Для фьючерса особенно важно, что фьючерсная цена уже является удобным объектом для риск-нейтральной оценки.

\subsubsection{Почему по MAE GARCH почти равен HV\_63d, а по TVNAE лучше?}

MAE считает ошибку по полной цене, где у deep ITM длинных опционов доминирует внутренняя стоимость. TVNAE делит ошибку на временную стоимость, то есть на компоненту, которую модель действительно должна объяснять. Поэтому TVNAE лучше выявляет преимущество GARCH в оценке опциональной части цены.

\subsubsection{Что означает «GARCH и HV\_63d статистически неразличимы»?}

Это не значит, что их численные MAE абсолютно одинаковы. Это значит, что на полной выборке разница их дневных функций потерь не достаточно велика относительно вариации, чтобы отвергнуть гипотезу равного качества.

\subsubsection{Почему regression $R^2=0{,}27$ не мал?}

Для финансовых микроданных ошибка зависит от множества ненаблюдаемых факторов: ликвидности, спроса и предложения, implied volatility surface, особенностей расчётной цены. Поэтому объяснить 27\% вариации простыми характеристиками контракта --- содержательный результат, а не провал.

\subsubsection{Почему регрессия не доказывает причинность?}

OLS здесь описывает связь между характеристиками и абсолютной ошибкой. Наблюдения не являются результатом эксперимента, поэтому нельзя сказать, что изменение DTE причинно создаёт ошибку. Корректная формулировка: DTE и денежность статистически связаны с масштабом ошибки.

\subsubsection{Почему ошибка выше на длинных опционах?}

Цена длинного опциона зависит от ожиданий волатильности на длительном горизонте. Одна текущая sigma плохо описывает всю будущую траекторию риска. Поэтому для длинных DTE нужна срочная структура волатильности, а не одно число.

\subsubsection{Почему в high vol GARCH не лучший, хотя он модель волатильности?}

GARCH хорошо описывает кластеризацию волатильности, но не обязательно идеально работает при смене режима. В стрессовом состоянии модель может слишком сильно переносить текущий шок в будущее или запаздывать. Кроме того, high-vol подвыборка отличается по DTE и денежности.

\subsubsection{Что означает «премия раннего исполнения»?}

Это дополнительная стоимость американского права исполнить опцион раньше экспирации. В работе она оценивается как разность между CRR и Black-76 при одинаковой HV\_63d. Если премия положительная, американский стиль исполнения действительно добавляет стоимость.

\subsubsection{Почему Black-76 иногда ниже внутренней стоимости?}

Black-76 оценивает европейский опцион, который нельзя исполнить немедленно. Поэтому его цена теоретически может оказаться ниже текущей внутренней стоимости. Для американского опциона это недопустимо, потому что держатель может сразу исполнить контракт.

\subsubsection{Какой практический вывод для выбора модели?}

Если оцениваются американские опционы на фьючерс, сначала нужно использовать модель с ранним исполнением, например CRR. Для коротких и среднесрочных контрактов предпочтительнее GARCH-волатильность. Для длинных контрактов HV\_63d является сильным и простым baseline. HV\_21d в этой выборке использовать как основной вариант нецелесообразно.

\subsubsection{Что главное улучшить в будущей работе?}

Главное улучшение --- построить поверхность подразумеваемой волатильности по сроку и денежности. Второй шаг --- добавить ликвидностные переменные: bid-ask spread, объёмы, открытый интерес. Третий --- проверить модели со стохастической волатильностью, jumps или regime-switching.

\section{Мини-шпаргалка с числами}

\begin{itemize}
    \item Период: 3 января 2024 --- 19 мая 2026.
    \item Исходная база: 19 574 наблюдения.
    \item После фильтрации: 19 565 наблюдений.
    \item Общая выборка для трёх CRR-моделей: 18 058 наблюдений.
    \item MAE: GARCH 107,0; HV\_63d 107,1; HV\_21d 119,3.
    \item RMSE: GARCH 145,8; HV\_63d 140,6; HV\_21d 156,4.
    \item Mean error: GARCH $-78{,}1$; HV\_63d $-89{,}7$; HV\_21d $-95{,}3$.
    \item TVNAE mean: GARCH 0,66; HV\_63d 0,70; HV\_21d 0,76.
    \item Black-76 + HV\_63d: MAE 156,2, то есть примерно на 46\% хуже CRR + HV\_63d.
    \item Средняя премия раннего исполнения: 59,28 пункта.
    \item Доля случаев, где Black-76 ниже внутренней стоимости: 24,5\%.
    \item Тест средней ошибки: все модели $\mu<0$, $p<0{,}001$.
    \item DM-test: GARCH и HV\_63d статистически неразличимы; HV\_21d значимо хуже HV\_63d, $p=0{,}0003$.
    \item OLS: $R^2\approx0{,}27$; DTE и денежность важнее простого индикатора high vol.
\end{itemize}

\section{Самый короткий финальный ответ}

Если после доклада попросят сформулировать результат в двух-трёх предложениях, можно сказать так:

\textit{В работе показано, что для опционов на фьючерс MXI американская постановка принципиально важна: Black-76 существенно занижает цены, потому что не учитывает раннее исполнение. Среди способов задания волатильности GARCH лучше описывает временную стоимость и короткие горизонты, но на полной выборке статистически близок к устойчивой 63-дневной исторической волатильности. Все простые модели систематически недооценивают рынок, поэтому дальнейшее развитие должно идти через поверхность подразумеваемой волатильности, ликвидность и более гибкие динамические модели.}

\end{document}
